% =====================================================================
%  APPENDIX: The Jaynes separation as a fibration --- the density
%  matrix as invariant ring, the missing global section, and the two
%  kinds of forgetting. Paper 1, "It from Bit via Goedel".
%  CREATED 2026-07-19. Expands the main-body insert
%  jaynes_paragraph_2026-07-19.tex; consolidates the session records
%  of 2026-07-19 and elevates the pure/mixed-state observations of the
%  2026-04-11 session (operational indistinguishability of partial
%  trace and Bayesian updating) to their theorem form (HJW).
%
%  USAGE. Plain \section --- no appendices wrapper in this file; wrap
%  all appendix inputs once in the main document, as for
%  appendix_ghz_w_proofs_2026-06-11.tex.
%
%  PREAMBLE. Requires: amsthm environments (proposition, remark ---
%  the starred remark style as in the GHZ/W appendix), \ket/\bra
%  (braket or physics package), biblatex \autocite.
%
%  CROSS-REFERENCES to the main text (adjust to master-file labels):
%    sec:hopf           -- the Hopf-fibration picture
%    sec:born           -- the degree-two Born argument
%    sec:fanout         -- the FANOUT logic-measure correspondence
%    sec:self-reference -- the diagonal obstruction
%    sec:contextuality  -- Hopf nontriviality = single-qubit
%                          contextuality (the framework's
%                          identification; see honesty flag in
%                          \S\ref{app:jaynes:section})
%  Companion-paper cross-references are marked with %% PAPER2 / %% PAPER3
%  comments; replace with \autocite of the companion papers when their
%  citation keys exist.
%
%  NEW CITATION KEYS used in this file (RIS entries for Mendeley
%  import in jaynes_references_2026-07-19.ris):
%    jaynes1990probability, schrodinger1936probability,
%    hughston1993ensemble, steenrod1951topology, berry1984quantal,
%    simon1983holonomy, wootters1981statistical, bengtsson2006geometry,
%    atiyah1982convexity, selinger2007dagger, coecke2017picturing,
%    abramsky2011sheaf, zurek2003decoherence, landauer1961irreversibility,
%    cleve1998quantum, wigner1961remarks, bargmann1954unitary,
%    pusey2012reality, harrigan2010einstein.
%  Several of these (jaynes, zurek, landauer, coecke2017, abramsky2011)
%  may already be in the Mendeley library under existing keys; check
%  before importing, and reconcile keys rather than duplicating.
%
%  STATUS DISCIPLINE. The working-document tags map onto
%  \S\ref{app:jaynes:status} as follows: STRUCTURAL = "established
%  mathematics"; CANDIDATE = "the framework's assembly"; RHYME =
%  "a fenced analogy". VERIFY flags: page ranges in the RIS file are
%  from memory and should be checked at final drafting.
% =====================================================================

\section{The Jaynes separation as a fibration: the density matrix, the
  missing section, and the two kinds of forgetting}
\label{app:jaynes}

Jaynes's complaint \autocite{jaynes1990probability} is that the
density-matrix formalism scrambles ``realities of Nature'' with
``incomplete human information'' and that no one has shown how to
unscramble them. The main text answers with a claim: both ingredients
are epistemic, but the ignorance differs in kind --- classical
probability is ignorance of an external state, quantum probability is
the observer's ignorance of their own state, carried by the global
phase. The immediate objection is that the formalism has already
discarded the phase: $\rho = \ket{\psi}\bra{\psi}$ is phase-invariant,
and every prediction factors through $\rho$. This appendix develops
the answer given in outline in the main text, in four steps: the
density matrix is the invariant ring of the phase action, not a
formalism from which the phase is absent
(\S\ref{app:jaynes:invariants}); the quotient leaves a quantum residue
because the bundle admits no global section, and what survives the
quotient is the fibre's curvature (\S\ref{app:jaynes:section}); for
mixed states the two kinds of probability enter through two distinct
forgetful operations whose provenance is recorded in the FANOUT
structure, not in the state (\S\ref{app:jaynes:mixed}); and the
unobservability of a phase is boundary-relative, with exactly one
phase --- that of the self-inclusive total --- unobservable for
structural rather than contingent reasons
(\S\ref{app:jaynes:boundary}). A closing section separates the
established mathematics from the framework's assembly, fences one
analogy, and answers the $\psi$-epistemic objection
(\S\ref{app:jaynes:status}). Throughout, ``the bundle'' means the
principal $U(1)$-bundle $S^{2n-1} \to \mathbb{C}P^{n-1}$ of
\S\ref{sec:hopf}, of which the qubit case $S^{3} \to S^{2}$ is the
Hopf fibration proper.

\subsection{The density matrix as the invariant ring of the phase
  action}
\label{app:jaynes:invariants}

Let $U(1)$ act on unit vectors $\psi \in \mathbb{C}^{n}$ by
$\psi \mapsto e^{i\alpha}\psi$. The framework's reading of this
action, established in \S\ref{sec:self-reference} and
\S\ref{sec:born}, is that the fibre coordinate is the observer's
self-referential ignorance: the one datum about the total
configuration that the observer, being part of it, cannot FANOUT. The
question is then what an observer so constituted can consistently
write down. The answer is: exactly the invariants.

\begin{proposition}[The density matrix is a complete invariant]
\label{prop:invariant-ring}
(i) Two unit vectors satisfy
$\ket{\psi}\bra{\psi} = \ket{\varphi}\bra{\varphi}$ if and only if
$\varphi = e^{i\alpha}\psi$ for some $\alpha$. (ii) Every
$U(1)$-invariant polynomial in the components $(\psi, \bar\psi)$ is a
polynomial in the entries $\psi_i\bar\psi_j$ of
$\ket{\psi}\bra{\psi}$, and these entries span the invariants of
minimal degree.
\end{proposition}

\begin{proof}
(i) A rank-one projector determines its range, hence the ray. (ii) A
monomial $\psi^{a}\bar\psi^{b}$ of bidegree $(|a|,|b|)$ picks up
$e^{i\alpha(|a|-|b|)}$ under the action, so invariance forces
$|a| = |b|$; there are consequently no invariants of bidegree $(1,0)$
or $(0,1)$, and the bidegree-$(1,1)$ invariants are spanned by the
$\psi_i\bar\psi_j$. Any monomial of bidegree $(k,k)$ is a product of
$k$ such factors.
\end{proof}

The proposition is elementary, but its reading is the point. The
density matrix is not a formalism in which the phase happens to be
absent; it is, by construction, \emph{everything except the phase} ---
the invariant ring of the self-ignorance symmetry, packaged as an
operator. Its diagonal entries are the Born weights of
\S\ref{sec:born}; its off-diagonal entries are the relative phases,
which is to say the interference data. Three consequences follow.

First, the Born rule is quadratic in $\psi$ but \emph{linear} in
$\rho$. There is no tension between these facts: the quadratic toll
identified in \S\ref{sec:born} --- the minimum-degree passage from
fibre coordinates to public quantities --- is paid exactly once, in
forming the invariants, after which every expectation value is linear.
The squaring \emph{is} the crossing.

Second, the shape of the object is the fingerprint of the quotient. A
classical epistemic state is a point of a simplex because classically
every quantity is in principle observable: the symmetry group of the
ignorance is trivial, and the invariant ring is the full coordinate
ring. A quantum epistemic state is a point of the invariant ring of a
nontrivial action, and invariant rings of nontrivial actions carry
relations --- here Hermiticity, positivity, and (for pure states)
rank one, with the off-diagonal structure as the residue of the
quotiented fibre. That quantum states are operators rather than
distributions is not an unexplained formal peculiarity; it is what the
invariant ring of a $U(1)$ action looks like.

Third, the construction has a categorical name. The passage
$\psi \mapsto \ket{\psi}\bra{\psi}$ is the doubling of categorical
quantum mechanics --- Selinger's CPM construction
\autocite{selinger2007dagger}, the doubled category of Coecke and
Kissinger \autocite{coecke2017picturing} --- which annihilates global
phases by construction and, augmented with discarding, generates
exactly the completely positive maps on density matrices. The
framework's dictionary of classical structures
(\S\ref{sec:fanout}) sits downstream of this doubling; the present
appendix supplies its epistemic reading.

\subsection{Why the quotient leaves a quantum residue: no section, but
  a connection}
\label{app:jaynes:section}

If the phase could be split off cleanly --- an objective ray together
with a subjective phase coordinate --- then the state space would be a
product, $S^{2n-1} \cong \mathbb{C}P^{n-1} \times S^{1}$, and
amplitudes would be classical probabilities wearing a decorative
label. The obstruction to this split is the content of the following
standard facts.

\begin{proposition}[No global phase convention]
\label{prop:no-section}
A principal bundle admits a global section if and only if it is
trivial \autocite{steenrod1951topology}. The bundle
$S^{2n-1} \to \mathbb{C}P^{n-1}$ is nontrivial for every $n \geq 2$:
its associated line bundle has first Chern class generating
$H^{2}(\mathbb{C}P^{n-1};\mathbb{Z}) \cong \mathbb{Z}$; equivalently,
for $n = 2$, the clutching map is the generator of
$\pi_{1}(U(1)) \cong \mathbb{Z}$, which is the statement that the Hopf
map has Hopf invariant one. Consequently there is no continuous
assignment of a representative state vector to every ray: no global
phase convention exists.
\end{proposition}

The interpretive step --- the framework's, not the literature's --- is
to read this as the precise formal content of Jaynes's
unscramblability. A separation of the description into ``objective
part'' and ``subjective part'' would be a global section: a consistent
rule, defined once for the whole state space, saying which datum is
the world and which is the observer's ignorance of themselves. Such a
rule exists chart by chart --- that is exactly what a local phase
convention, a gauge, is --- and Proposition~\ref{prop:no-section} says
the charts cannot be glued. Classical ignorance is a product, world
$\times$ credence, and a product unscrambles by projection.
Self-referential ignorance is a twisted bundle, and the twist is
measured by a nonvanishing class. The omelette does not unscramble
pointwise; it \emph{classifies}.

This does not mean the quotient discards the fibre's structure. What
survives is the fibre's differential geometry: the natural connection
on the associated line bundle, whose holonomy is the geometric phase
--- observable entirely within the projective description
\autocite{berry1984quantal, simon1983holonomy} --- and whose curvature
is the Fubini--Study form. The real part of the Fubini--Study
structure is the metric of statistical distinguishability
\autocite{wootters1981statistical}; the imaginary part is the
symplectic form whose Hamiltonian flow is the Schr\"odinger equation
\autocite{bengtsson2006geometry}. In one sentence: the observer
quotients out the phase but keeps its curvature. On this reading
Jaynes's two ingredients are assigned as follows --- the ``realities
of Nature'' are the curvature, chart-invariant and public; the
``incomplete human information'' is the fibre coordinate, chartful and
private, together with the classical mixing weights of
\S\ref{app:jaynes:mixed}.

\begin{remark}[Relation to contextuality, with an honesty flag]
Abramsky and Brandenburger define contextuality as the nonexistence of
a global section of an empirical presheaf \autocite{abramsky2011sheaf}.
Proposition~\ref{prop:no-section} is the continuous, bundle-theoretic
counterpart of that schema, and the main text's identification of Hopf
nontriviality with single-qubit contextuality
(\S\ref{sec:contextuality}) is the bridge between the two. The flag:
standard Kochen--Specker contextuality requires Hilbert-space
dimension at least three and is silent about a single qubit. The
single-qubit usage here is the framework's bundle-theoretic extension
of the term, argued in the main text; to our knowledge the discrete
sheaf-theoretic and continuous bundle-theoretic obstruction
literatures have not been joined in print, and we do not attribute the
identification to either.
\end{remark}

\subsection{Mixed states: two forgetful operations, and where the
  provenance lives}
\label{app:jaynes:mixed}

The general density matrix
$\rho = \sum_i p_i \ket{\psi_i}\bra{\psi_i}$ is built by two
operations, and the two kinds of probability in Jaynes's omelette
enter one through each. The first is the projective quotient of
\S\ref{app:jaynes:invariants}: its kernel is the fibre coordinate, the
observer's ignorance of their own state, unresolvable for the
structural reason of \S\ref{sec:self-reference}. The second is the
partial trace together with convex mixing: its kernel is correlation
with systems outside the description --- by purification, every mixed
$\rho_A$ is $\mathrm{Tr}_B\ket{\Psi}\bra{\Psi}$ for a pure state on a
larger system --- and this ignorance is resolvable in principle,
because whoever holds $B$, or a record of the preparation, can supply
the missing fact. The framework's founding pair of sentences ---
quantum probability is ignorance of one's own state, classical
probability is ignorance of an external state --- is thereby realised
as the pair of kernels of the two maps that construct $\rho$.

The scrambling is then a theorem rather than a confusion. Schr\"odinger
observed the ensemble ambiguity \autocite{schrodinger1936probability};
Hughston, Jozsa and Wootters completed it: every ensemble
decomposition of a given $\rho$ is realisable by a suitable
measurement on a purifying system \autocite{hughston1993ensemble}.
The two mechanisms --- tracing out entanglement, and mixing over
preparation ignorance --- are therefore operationally
indistinguishable at the level of $\rho$ itself. In statistical
language: $\rho$ is a sufficient statistic for all future predictions,
and sufficiency is forgetting. Jaynes demanded that the separation be
exhibited within the state object; the theorem says the state object
is precisely the invariant that forgets it. Seventy-five years of
failure to unscramble the omelette inside the formalism is thereby
explained rather than lamented.

Where, then, does the separation live? In the provenance. Two
decompositions of the same $\rho$ correspond to two candidate
histories of irreversible FANOUT events --- which copies were actually
made, by whom, of what --- and the actual history disambiguates. The
correspondence of \S\ref{sec:fanout} supplies the formal home for that
record structure. The slogan form: the omelette is unscrambled in the
circuit, not in the state. We flag honestly that the composite
statement --- FANOUT record structure plus the
Hughston--Jozsa--Wootters ambiguity yields a well-defined
classical/quantum provenance for any $\rho$ arising in a given circuit
--- is here a definition plus two theorems standing adjacent to one
another; we have not written the composite as a single theorem, and
\S\ref{app:jaynes:status} records it as the natural point of attack
for a critical reader.

The conversion between the two kinds of ignorance also has a clean
geometric expression. Complete dephasing in a pointer basis
\autocite{zurek2003decoherence} maps $\rho$ to its diagonal; restricted
to pure states it is the map
\[
  \mathbb{C}P^{n-1} \longrightarrow \Delta^{n-1}, \qquad
  [\psi] \longmapsto \bigl(|\psi_1|^2, \dots, |\psi_n|^2\bigr),
\]
which is the moment map of the residual torus action on
$\mathbb{C}P^{n-1}$ --- the relative phases that survived the first
quotient --- and its image is the probability simplex by the convexity
theorem \autocite{atiyah1982convexity}. The chain
\[
  S^{2n-1} \;\longrightarrow\; \mathbb{C}P^{n-1}
  \;\longrightarrow\; \Delta^{n-1}
\]
thus separates the two quotients in kind: the first is private and
structural, effected by the observer's constitution; the second is
public and processual, effected by FANOUT in a pointer basis,
converting fibre-type, interference-capable ignorance into
simplex-type, record-resolvable ignorance. The second quotient is
irreversible and carries Landauer's toll
\autocite{landauer1961irreversibility}; its thermodynamic ledger is
the subject of the companion paper.
%% PAPER3: replace the preceding clause's prose pointer with a
%% \autocite of Paper 3 when its key exists.

\subsection{Two grades of unobservability: ``global'' is
  boundary-relative}
\label{app:jaynes:boundary}

One residual form of the objection remains: if the density matrix
``eliminates'' the phase, in what sense is the phase still anything at
all? The answer is that a phase is global only relative to a system
boundary, and the boundary can usually be moved. Phase kickback makes
this operational \autocite{cleve1998quantum}: a controlled unitary
acting on an eigenstate of its target writes the target's eigenphase
--- global, hence unobservable, on the target alone --- onto the
relative phase of the control, where it interferes. Wigner's friend is
the same fact with the observer in the target's role
\autocite{wigner1961remarks}: the friend's overall phase is
unobservable to the friend and interference-accessible to Wigner. A
subsystem's global phase is therefore only \emph{contingently}
unobservable: enlarge the boundary and it becomes relative.

Exactly one phase resists this manoeuvre: that of the self-inclusive
total, because the observer cannot enlarge the boundary past
themselves. That is the diagonal obstruction of
\S\ref{sec:self-reference}, and that phase --- not the movable
subsystem phases --- is what the framework means by the observer's
self-ignorance. There are thus two grades of unobservability,
contingent and structural, and only the second is self-ignorance
proper. The density matrix of a subsystem removes the contingent grade
reversibly --- an outside party can restore the phase by interference
--- while the projective quotient on the self-inclusive description
removes the structural grade from every chart, permanently. Even then
the curvature shows through, as \S\ref{app:jaynes:section} records:
interference is the structural grade's public shadow. The tower of
observers that this picture generates, and its termination, are taken
up in the companion paper on measurement-scale unification.
%% PAPER2: replace the preceding clause's prose pointer with a
%% \autocite of Paper 2 when its key exists.

\subsection{Status of the claims, one fenced analogy, and the
  \texorpdfstring{$\psi$}{psi}-epistemic objection}
\label{app:jaynes:status}

\emph{Established mathematics.} Proposition~\ref{prop:invariant-ring}
is elementary invariant theory; Proposition~\ref{prop:no-section} is
Steenrod plus the computation of a Chern class; the geometric phase as
holonomy and the Fubini--Study structure as metric-plus-symplectic
form are standard \autocite{berry1984quantal, simon1983holonomy,
wootters1981statistical, bengtsson2006geometry}; the ensemble
ambiguity is Schr\"odinger and Hughston--Jozsa--Wootters
\autocite{schrodinger1936probability, hughston1993ensemble}; the
moment-map identification of the simplex is the convexity theorem
\autocite{atiyah1982convexity}; the doubling construction is Selinger
and Coecke--Kissinger \autocite{selinger2007dagger,
coecke2017picturing}. None of this is ours, and none of it is at
stake.

\emph{The framework's assembly.} Three readings are ours and stand or
fall with the main text's derivation chain: that the fibre coordinate
is the formal carrier of self-referential ignorance
(\S\ref{sec:self-reference}, \S\ref{sec:born}); that the nonvanishing
obstruction class of Proposition~\ref{prop:no-section} is the formal
content of Jaynes's unscramblability; and that the classical/quantum
provenance of a given $\rho$ is fixed by the FANOUT record rather than
by the state. The first two inherit whatever support the main text's
chain has earned. The third is, as flagged in
\S\ref{app:jaynes:mixed}, a definition and two theorems not yet
composed into one; a reader wishing to break the framework's answer to
Jaynes should push exactly there, and a reader wishing to complete it
should prove exactly that.

\emph{A fenced analogy.} Quotienting phases at the level of states
forces central extensions at the level of symmetry groups: projective
unitary representations of the Galilei group are true representations
of its central extension, and the central charge is the mass, with
Bargmann's superselection rule as enforcement
\autocite{bargmann1954unitary}. The phase one cannot see returns in
the group cohomology carrying a mass label --- an arresting echo of
the framework's identification of mass with un-narratable debt. The
fence: the mechanism is Galilei-specific; the corresponding extensions
of the Poincar\'e group are trivial and mass enters through the
Casimir instead. We record the rhyme and build nothing on it.

\emph{The $\psi$-epistemic objection.} A reader who hears ``quantum
probability is epistemic'' will reach for the theorem of Pusey,
Barrett and Rudolph \autocite{pusey2012reality}, which --- within the
ontological-models framework of Harrigan and Spekkens
\autocite{harrigan2010einstein}, and granted preparation independence
--- excludes models in which distinct pure states correspond to
overlapping probability distributions over an observer-independent
ontic state space $\Lambda$. The present reading is not such a model,
and not by evasion: the ignorance asserted here is reflexive, and
there is no observer-independent $\Lambda$ over which it could be a
distribution --- Proposition~\ref{prop:no-section} is the geometric
expression of the absence of any global assignment of ``the underlying
value''. An ontological model in the Harrigan--Spekkens sense is,
structurally, a section-like object, and the bundle has no section. We
state the residue honestly: this is a structural observation, not a
theorem that the framework evades the Pusey--Barrett--Rudolph
constraint, because we have not formalised the framework as, or
against, an ontological model. A careful placement of the framework
relative to the Harrigan--Spekkens taxonomy is an open item, and we
would welcome it being carried out by hands other than ours.

% =====================================================================
% BIBLIOGRAPHY STUBS (informal; authoritative RIS entries in
% jaynes_references_2026-07-19.ris; VERIFY page ranges at final
% drafting)
% jaynes1990probability: E. T. Jaynes, Probability in quantum theory,
%   in Complexity, Entropy, and the Physics of Information (ed. W. H.
%   Zurek), Addison-Wesley (1990), 381-403.
% schrodinger1936probability: E. Schroedinger, Probability relations
%   between separated systems, Math. Proc. Camb. Phil. Soc. 32 (1936),
%   446-452.
% hughston1993ensemble: L. P. Hughston, R. Jozsa, W. K. Wootters, A
%   complete classification of quantum ensembles having a given
%   density matrix, Phys. Lett. A 183 (1993), 14-18.
% steenrod1951topology: N. Steenrod, The Topology of Fibre Bundles,
%   Princeton University Press (1951).
% berry1984quantal: M. V. Berry, Quantal phase factors accompanying
%   adiabatic changes, Proc. R. Soc. Lond. A 392 (1984), 45-57.
% simon1983holonomy: B. Simon, Holonomy, the quantum adiabatic
%   theorem, and Berry's phase, Phys. Rev. Lett. 51 (1983), 2167-2170.
% wootters1981statistical: W. K. Wootters, Statistical distance and
%   Hilbert space, Phys. Rev. D 23 (1981), 357-362.
% bengtsson2006geometry: I. Bengtsson, K. Zyczkowski, Geometry of
%   Quantum States, Cambridge University Press (2006).
% atiyah1982convexity: M. F. Atiyah, Convexity and commuting
%   Hamiltonians, Bull. London Math. Soc. 14 (1982), 1-15.
% selinger2007dagger: P. Selinger, Dagger compact closed categories
%   and completely positive maps, ENTCS 170 (2007), 139-163.
% coecke2017picturing: B. Coecke, A. Kissinger, Picturing Quantum
%   Processes, Cambridge University Press (2017).
% abramsky2011sheaf: S. Abramsky, A. Brandenburger, The sheaf-theoretic
%   structure of non-locality and contextuality, New J. Phys. 13
%   (2011), 113036.
% zurek2003decoherence: W. H. Zurek, Decoherence, einselection, and
%   the quantum origins of the classical, Rev. Mod. Phys. 75 (2003),
%   715-775.
% landauer1961irreversibility: R. Landauer, Irreversibility and heat
%   generation in the computing process, IBM J. Res. Dev. 5 (1961),
%   183-191.
% cleve1998quantum: R. Cleve, A. Ekert, C. Macchiavello, M. Mosca,
%   Quantum algorithms revisited, Proc. R. Soc. Lond. A 454 (1998),
%   339-354.
% wigner1961remarks: E. P. Wigner, Remarks on the mind-body question,
%   in The Scientist Speculates (ed. I. J. Good), Heinemann (1961),
%   284-302.
% bargmann1954unitary: V. Bargmann, On unitary ray representations of
%   continuous groups, Ann. Math. 59 (1954), 1-46.
% pusey2012reality: M. F. Pusey, J. Barrett, T. Rudolph, On the
%   reality of the quantum state, Nature Physics 8 (2012), 475-478.
% harrigan2010einstein: N. Harrigan, R. W. Spekkens, Einstein,
%   incompleteness, and the epistemic view of quantum states, Found.
%   Phys. 40 (2010), 125-157.
% =====================================================================
